% ==============================================================================
% Section 4: Sector-wise FDI Analysis
% ==============================================================================

\section{Sector-wise FDI Analysis}

Decomposing direct investment capital stocks by economic activity reveals stark structural contrasts between the composition of foreign capital invested inside Switzerland and the overseas capital footprint of Swiss enterprises abroad. Switzerland functions predominantly as a specialized service, trading, and holding platform for inward investors, whereas Swiss outward direct investment maintains an exceptionally robust high-tech manufacturing core.

\subsection{Inward FDI Allocation Across Sectors}

Inward foreign direct investment in Switzerland is heavily concentrated in the tertiary (services) sector, which accounts for more than three-quarters of total stock, as detailed in Table~\ref{tab:inward_sector}.

\begin{table}[H]
\centering
\small
\caption{Sectoral Composition of Inward FDI Capital Stocks in Switzerland (2024)}
\label{tab:inward_sector}
\begin{tabularx}{\textwidth}{>{\raggedright}p{5.2cm} r r >{\raggedright\arraybackslash}X}
\toprule
\rowcolor{tableheadbg}
\textbf{Economic Sector / Activity} & \textbf{Capital Stock ($\CHF$\,mn)} & \textbf{Share (\%)} & \textbf{Primary Industrial Drivers} \\
\midrule
Chemicals and Plastics & 109,860 & 11.9\% & Specialty chemicals, active ingredients \\
Electronics, Energy, Optical & 52,679 & 5.7\% & Precision medical devices, instrumentation \\
Metals and Machinery & 24,894 & 2.7\% & Advanced tooling, automation robotics \\
Other Manufacturing & 20,338 & 2.2\% & Food processing, construction materials \\
\midrule
\rowcolor{tablerowbg}
\textbf{Total Manufacturing} & \textbf{207,770} & \textbf{22.4\%} & \textit{High-value, R\&D-intensive production} \\
\midrule
Trade (Wholesale \& Commodity) & 235,172 & 25.4\% & Global commodity hubs (Geneva/Zug) \\
Other Services (Holding \& Admin) & 443,374 & 47.9\% & Corporate headquarters, IP management \\
Other Services (General) & 52,294 & 5.6\% & IT services, consulting, logistics \\
Commercial Banks & 21,326 & 2.3\% & Wealth management, custody operations \\
Insurance Companies & 18,068 & 2.0\% & Reinsurance, life \& non-life underwriting \\
\midrule
\rowcolor{tablerowbg}
\textbf{Total Services} & \textbf{717,940} & \textbf{77.6\%} & \textit{Trade finance, headquarters \& services} \\
\midrule
\textbf{Total All Sectors} & \textbf{925,711} & \textbf{100.0\%} & \textbf{Complete Inward FDI Stock} \\
\bottomrule
\end{tabularx}
\vspace{0.3em}
\footnotesize
\textit{Source: Swiss National Bank (SNB). Figures in millions of Swiss Francs ($\CHF\,$mn) at year-end 2024.}
\normalsize
\end{table}

\noindent\textbf{Key Insights on Inward Sectoral Distribution:}
\begin{itemize}[leftmargin=1.5em, itemsep=2.5pt]
    \item \textbf{Dominance of the Tertiary Sector ($77.6\pct$):} Inward investors utilize Switzerland primarily for high-level tertiary operations. The single largest operational sector is \textbf{Wholesale Trade} ($\CHF\,235.2\bn$, $25.4\pct$), driven by Switzerland's status as the global capital for physical commodity trading (crude oil, metals, agricultural commodities clustered in Geneva, Zug, and Lugano).
    \item \textbf{Corporate Headquarters and Holding Entities:} Over $47\pct$ of total inward capital resides in holding, administrative, and corporate service vehicles, reflecting Switzerland's attractive IP holding regimes and bilateral treaty protections.
    \item \textbf{Manufacturing Anchored in Life Sciences:} Within secondary industry ($22.4\pct$), the \textbf{Chemicals and Plastics} segment represents over half of all manufacturing investment ($\CHF\,109.9\bn$), underscoring the pull of the Swiss life-sciences ecosystem.
\end{itemize}

\subsection{Outward FDI Allocation Across Sectors}

Swiss outward direct investment displays a markedly different structural profile. While services continue to absorb the majority of capital, Swiss industry maintains an extraordinarily high foreign manufacturing commitment of \textbf{$37.9\pct$} (Table~\ref{tab:outward_sector}).

\begin{table}[H]
\centering
\small
\caption{Sectoral Composition of Swiss Outward FDI Capital Stocks Abroad (2024)}
\label{tab:outward_sector}
\begin{tabularx}{\textwidth}{>{\raggedright}p{5.5cm} r r >{\raggedright\arraybackslash}X}
\toprule
\rowcolor{tableheadbg}
\textbf{Economic Sector / Activity} & \textbf{Capital Stock ($\CHF$\,mn)} & \textbf{Share (\%)} & \textbf{Swiss Multinational Anchors} \\
\midrule
Chemicals and Plastics & 190,217 & 14.2\% & Novartis, Roche, Lonza, Givaudan \\
Other Manufacturing (Food/Consumer) & 169,670 & 12.7\% & Nestlé, Lindt \& Sprüngli \\
Electronics, Energy, Optical & 86,864 & 6.5\% & ABB, TE Connectivity, Swatch Group \\
Metals and Machinery & 61,249 & 4.6\% & Schindler, Sulzer, Georg Fischer \\
\midrule
\rowcolor{tablerowbg}
\textbf{Total Manufacturing} & \textbf{507,999} & \textbf{37.9\%} & \textit{Specialized advanced industrial base} \\
\midrule
Trade & 187,146 & 14.0\% & Global procurement, distribution hubs \\
Insurance Companies & 109,810 & 8.2\% & Swiss Re, Zurich Insurance Group \\
Commercial Banks & 65,649 & 4.9\% & UBS Group (post-Credit Suisse merger) \\
Other Services \& Holding & 469,537 & 35.0\% & Regional management \& treasury hubs \\
\midrule
\rowcolor{tablerowbg}
\textbf{Total Services} & \textbf{832,142} & \textbf{62.1\%} & \textit{Financial services, reinsurance \& logistics} \\
\midrule
\textbf{Total All Sectors} & \textbf{1,340,141} & \textbf{100.0\%} & \textbf{Complete Outward FDI Stock} \\
\bottomrule
\end{tabularx}
\vspace{0.3em}
\footnotesize
\textit{Source: Swiss National Bank (SNB). Figures in millions of Swiss Francs ($\CHF\,$mn) at year-end 2024.}
\normalsize
\end{table}

% ------------------------------------------------------------------------------
% Pgfplots Figure 4.1: Comparative Sectoral Distribution (Inward vs Outward)
% Greyscale Bar Chart with Pattern Shading
% ------------------------------------------------------------------------------
\begin{figure}[H]
\centering
\begin{tikzpicture}
\begin{axis}[
    ybar=3pt,
    bar width=0.55cm,
    width=13.5cm,
    height=6.8cm,
    grid=major,
    grid style={line width=0.2pt, color=midgray!30},
    ylabel={\small\textbf{Share of Total Capital Stock (\%)}},
    symbolic x coords={Manufacturing, Wholesale Trade, Finance \& Ins., Other Services},
    xtick=data,
    xticklabel style={font=\footnotesize\bfseries},
    ytick={0, 10, 20, 30, 40, 50, 60},
    ymin=0, ymax=60,
    legend style={
        at={(0.5,1.08)},
        anchor=south,
        legend columns=2,
        font=\footnotesize,
        draw=pureblack,
        line width=0.6pt
    },
    enlarge x limits=0.18
]

    % Inward FDI Bars (Dark Gray Fill with Solid Edge)
    \addplot[fill=darkgraytext!70!white, draw=pureblack, line width=0.8pt] 
        coordinates {(Manufacturing, 22.4) (Wholesale Trade, 25.4) (Finance \& Ins., 4.3) (Other Services, 47.9)};
    \addlegendentry{Inward FDI in Switzerland}

    % Outward FDI Bars (Pure Black Outline with Light Fill)
    \addplot[fill=boxfill, draw=pureblack, line width=1.1pt, postaction={pattern=north east lines, pattern color=midgray}] 
        coordinates {(Manufacturing, 37.9) (Wholesale Trade, 14.0) (Finance \& Ins., 13.1) (Other Services, 35.0)};
    \addlegendentry{Swiss Outward FDI Abroad}

\end{axis}
\end{tikzpicture}
\caption{Cross-Directional Sectoral Comparison of Swiss Direct Investment Stocks (2024). Demonstrates the structural tilt of inward FDI toward trade/holding services and outward FDI toward advanced manufacturing and global reinsurance.}
\label{fig:sector_comparison}
\end{figure}

\subsection{Pharmaceutical and Chemical Sector Deep Dive}

The pharmaceutical, biotechnology, and specialty chemical cluster represents the indisputable crown jewel of Switzerland's direct investment landscape. 

\begin{itemize}[leftmargin=1.5em, itemsep=3pt]
    \item \textbf{Inward Trajectory ($\CHF\,109.9\bn$ in 2024):} Foreign direct investment in Swiss chemicals and pharmaceuticals expanded from $\CHF\,63.7\bn$ in 2015 to $\CHF\,109.9\bn$ in 2024, achieving a spectacular \textbf{$+72.5\pct$ cumulative growth rate}. Foreign healthcare leaders (such as Johnson \& Johnson, Merck, Takeda, and Pfizer) maintain critical active pharmaceutical ingredient (API) synthesis, biologics manufacturing, and European clinical development centers in Switzerland.
    \item \textbf{Outward Expansion ($\CHF\,190.2\bn$ in 2024):} Swiss pharmaceutical giants deployed $\CHF\,190.2\bn$ abroad in 2024 (up from $\CHF\,164.9\bn$ in 2015), representing $14.2\pct$ of Switzerland's total global direct investment assets. This overseas footprint reflects aggressive acquisitions of US biotech start-ups, oncology platforms, and global distribution entities.
\end{itemize}

\begin{mnecase}{The Basel Life-Sciences Triad: Roche, Novartis, and Inward Biopharma}{basel_cluster}
The canton of Basel-Stadt and surrounding Basel-Landschaft represent the world's densest life-sciences cluster. Together, \textbf{Novartis AG} and \textbf{F. Hoffmann-La Roche AG} deploy tens of billions in global R\&D, driving reciprocal cross-border investment:
\begin{itemize}[leftmargin=1.2em, itemsep=2pt]
    \item \textbf{Outward Pipeline Integration:} Both firms execute strategic overseas cross-border acquisitions (e.g., Roche's acquisition of Genentech and Spark Therapeutics; Novartis' acquisitions of AveXis and The Medicines Company) to secure advanced gene-therapy platforms in the United States.
    \item \textbf{Inward Synergies:} The clustering of world-class research institutes (Friedrich Miescher Institute, Biozentrum) attracts foreign biopharma suppliers and venture capital into Basel, creating a high-wage, patent-intensive agglomeration that is virtually impervious to conventional cost competition.
\end{itemize}
\end{mnecase}
